בפרק זה יפורטו כלי התוכנה, הספריות ותשתיות הפיתוח ששימשו לבניית
המערכת. הבחירה בכלים אלו נעשתה מתוך שיקולים הנדסיים של ביצועים בזמן
אמת, קלות תחזוקה ויכולת פריסה כשרת מרכזי.

\subsection{סביבות עבודה ותשתית חומרה}

בשל הבדלי המשאבים הנדרשים בין
שלב האימון לשלב ההרצה, נעשה שימוש בשתי סביבות נפרדות:
\begin{itemize}
\item \textbf{סביבת אימון (\textenglish{Cloud GPU}):} אימון מודלי
  ה־\textenglish{YOLO} בוצע על גבי פלטפורמת \textbf{\textenglish{Google
      Colab}} תוך שימוש במעבד גרפי (\textenglish{GPU}) \textenglish{T4} של חברת
  \textenglish{NVIDIA}. שימוש זה אפשר לקצר משמעותית את זמן האימון של
  הרשתות העמוקות.
\item \textbf{סביבת הרצה ובדיקה (\textenglish{Local CPU}):} בדיקת
  המערכת והרצת השרת בוצעו על מחשב מקומי בעל מעבד \textenglish{Intel
    i7-10700}. סביבה זו מדמה את תנאי העבודה האמיתיים בתחרות, בהם לא תמיד
  קיימת גישה למשאבי ענן חזקים.
\end{itemize}

\subsection{שפות תכנות, \textenglish{Frameworks} וניהול גרסאות}

כדי לעמוד בדרישות הפרויקט מבחינת ביצועים, ניידות ותחזוקה, נעשה שימוש
בכמה שכבות פיתוח שונות:
\begin{itemize}
\item \textbf{שפות תכנות:} צד השרת, עיבוד הוידאו והרצת המודלים פותחו
  ב־\textenglish{Python}, בעוד שממשק המשתמש פותח ב־\textenglish{Dart}
  באמצעות \textenglish{Flutter}.
\item \textbf{\textenglish{Frameworks}:} מודלי הזיהוי הורצו באמצעות
  \textenglish{Ultralytics} על גבי \textenglish{PyTorch}, השרת מומש
  באמצעות \textenglish{FastAPI} והאפליקציה נבנתה במסגרת
  \textenglish{Flutter}.
\item \textbf{ניהול גרסאות:} לאורך הפיתוח נעשה שימוש ב־\textenglish{Git}
  לניהול הקוד, מעקב אחרי שינויים וחזרה לגרסאות קודמות במקרה הצורך.
\end{itemize}

\subsection{ספריות ליבה ועיבוד נתונים (\textenglish{Backend})}

המערכת
מבוססת על שפת \textenglish{Python} ומשתמשת בספריות הבאות:
\begin{itemize}
\item \textbf{\textenglish{Ultralytics}}: המעטפת המרכזית להרצת
  מודלי ה־\textenglish{YOLO} בפורמט \textenglish{PyTorch}.
\item \textbf{\textenglish{Levenshtein}}: ספרייה לחישוב מהיר של
  מרחק עריכה בין מחרוזות. היא משמשת כמרכיב קריטי בתיקון שגיאות זיהוי של
  מספר הקבוצה על ידי השוואת הפלט למספרי הקבוצות המשתתפות במקצה.
\item \textbf{\textenglish{OpenCV} ו־\textenglish{SciPy}}: לניהול
  זרם הוידאו, גזירת תמונות (\textenglish{Crops}) והפעלת מסננים סטטיסטיים
  להחלקת מסלולי הרובוטים.
\item \textbf{\textenglish{NumPy}}: לניהול המטריצות של התמונות
  וביצוע חישובים מספריים על מיקומי האובייקטים.
\end{itemize}

\subsection{תשתית השרת (\textenglish{Infrastructure})}

כדי לאפשר עיבוד
וידאו כבד בסביבת אינטרנט, נעשה שימוש בכלים המאפשרים מקביליות וביצועים
גבוהים:
\begin{itemize}
\item \textbf{\textenglish{FastAPI} ו־\textenglish{Uvicorn}}: השרת
  מבוסס על \textenglish{FastAPI} ורץ על גבי מנוע
  \textenglish{Uvicorn}. בחירה זו מאפשרת טיפול מהיר בבקשות
  \textenglish{HTTP} והעלאת קבצים בצורה יעילה.
\item \textbf{\textenglish{asyncio}}: שימוש בתכנות אסינכרוני
  (\textenglish{Asynchronous Programming}) המאפשר לשרת לנהל מספר פעולות
  במקביל. תשתית זו קריטית כדי לאפשר למשתמשים לבדוק סטטוס עיבוד בזמן
  שהמעבד עסוק בחישוב המודלים.
\end{itemize}

\subsection{ממשק המשתמש (\textenglish{Frontend})}

האפליקציה פותחה
ב־\textbf{\textenglish{Flutter}} (שפת \textenglish{Dart}) תוך שימוש
בחבילות ייעודיות המבטיחות חוויית משתמש חלקה:
\begin{itemize}
\item \textbf{\textenglish{dio}}: לניהול בקשות
  ה־\textenglish{HTTP} והעברת סרטוני הוידאו מול השרת.
\item \textbf{\textenglish{media\_kit}}: חבילה מתקדמת להצגת וידאו
  וניגון מדיה ביציבות גבוהה על גבי מכשירים שונים.
\item \textbf{\textenglish{file\_picker}}: המאפשרת למשתמש לבחור את
  סרטוני המקצים מתוך זכרון המכשיר בצורה נוחה.
\end{itemize}

\subsection{תיוג נתונים וניהול מודלים}

לצורך יצירת מערך הנתונים, נעשה
שימוש בפלטפורמת \textbf{\textenglish{Roboflow}}. הכלי שימש לתיוג ידני
של הספרות, ביצוע אוגמנטציה (\textenglish{Augmentation}) וייצוא
הנתונים בפורמטים המתאימים לאימון.

\subsection{ריכוז הכלים וסביבות הפיתוח}
בשל ריבוי הספריות והכלים בהם נעשה שימוש, ריכזתי את הפירוט הטכני המלא בטבלה \ref{tab:tools_full} המופיעה בנספח \ref{app:tools}.

